% R586: EGA I, 0, 3.1.4--3.1.6.
% Traducción: OpenAI Codex, GPT-6.1 Sol, esfuerzo Ultra.
\begin{env}[3.1.4]
\label{0.3.1.4-fr}
Cuando $\Sigma$ es la especie de estructura de grupo o de anillo, el hecho de que
el prehaz $U\to\mathcal{F}(U)$ con valores en $\boldsymbol{K}$ sea un
haz de \emph{conjuntos} implica \emph{ipso facto} que es un
haz \emph{con valores en $\boldsymbol{K}$} (es decir, un haz de
grupos o de anillos en el sentido de (G))%
\footnote{Esto se debe a que, en la categoría $\boldsymbol{K}$, todo morfismo
que es una \emph{biyección} (como aplicación de conjuntos) es un
\emph{isomorfismo}. Esto ya no es cierto cuando $\boldsymbol{K}$ es, por ejemplo,
la categoría de los espacios topológicos.}. Pero esto deja de ser así
cuando, por ejemplo, $\boldsymbol{K}$ es la categoría de los \emph{anillos
topológicos} (con las representaciones continuas como morfismos): un
haz con valores en $\boldsymbol{K}$ es un haz de anillos
$U\to\mathcal{F}(U)$ tal que, para todo abierto $U$ y todo recubrimiento de $U$
por abiertos $U_\alpha\subset U$, la topología del anillo $\mathcal{F}(U)$
sea \emph{la menos fina} que hace continuas las representaciones
$\mathcal{F}(U)\to\mathcal{F}(U_\alpha)$. En este caso se dirá que
$U\to\mathcal{F}(U)$, considerado como haz de anillos (sin topología), es
\emph{subyacente} al haz de anillos topológicos
$U\to\mathcal{F}(U)$. Por tanto, los morfismos
$u_V:\mathcal{F}(V)\to\mathcal{G}(V)$ ($V$ abierto arbitrario de $X$) de
haces de anillos topológicos son homomorfismos de los haces de
anillos subyacentes tales que $u_V$ es \emph{continuo} para todo abierto
$V\subset X$; para distinguirlos de los homomorfismos cualesquiera de los
haces de anillos subyacentes, se los llamará homomorfismos
\emph{continuos} de haces de anillos topológicos. Se tienen definiciones y
convenciones análogas para los haces de espacios topológicos o de grupos
topológicos.
\end{env}

\oldpage[0\textsubscript{I}]{25}

\begin{env}[3.1.5]
\label{0.3.1.5-fr}
Es claro que, para toda categoría $\boldsymbol{K}$, si $\mathcal{F}$ es un
prehaz (resp. un haz) sobre $X$ con valores en $\boldsymbol{K}$ y $U$ es un
abierto de $X$, los $\mathcal{F}(V)$ para los abiertos $V\subset U$ constituyen
un prehaz (resp. un haz) con valores en $\boldsymbol{K}$, que se
llama prehaz (resp. haz) \emph{inducido} por $\mathcal{F}$ sobre $U$ y
se denota por $\mathcal{F}|U$.

Para todo morfismo $u:\mathcal{F}\to\mathcal{G}$ de prehaces sobre $X$ con
valores en $\boldsymbol{K}$, se denotará por $u|U$ el morfismo
$\mathcal{F}|U\to\mathcal{G}|U$ formado por los $u_V$ para $V\subset U$.
\end{env}

\begin{env}[3.1.6]
\label{0.3.1.6-fr}
Supongamos ahora que la categoría $\boldsymbol{K}$ admite \emph{límites
inductivos} (T, 1.8); entonces, para todo prehaz (y, en particular, todo
haz) $\mathcal{F}$ sobre $X$ con valores en $\boldsymbol{K}$ y todo $x\in X$,
se puede definir el \emph{tallo} $\mathcal{F}_x$ como el objeto de $\boldsymbol{K}$
que es el límite inductivo de los $\mathcal{F}(U)$ según el conjunto filtrante (para
$\supset$) de las vecindades abiertas $U$ de $x$ en $X$, y para los morfismos
$\rho_U^V:\mathcal{F}(V)\to\mathcal{F}(U)$. Si
$u:\mathcal{F}\to\mathcal{G}$ es un morfismo de prehaces con valores en
$\boldsymbol{K}$, se define, para todo $x\in X$, el morfismo
$u_x:\mathcal{F}_x\to\mathcal{G}_x$ como el límite inductivo de los
$u_U:\mathcal{F}(U)\to\mathcal{G}(U)$ según el conjunto de las vecindades abiertas
de $x$; así se define $\mathcal{F}_x$ como funtor covariante en
$\mathcal{F}$, con valores en $\boldsymbol{K}$, para todo $x\in X$.

Cuando $\boldsymbol{K}$ está además definida por una especie de estructura con
morfismos $\Sigma$, se siguen llamando \emph{secciones sobre $U$} de un
\emph{haz $\mathcal{F}$} con valores en $\boldsymbol{K}$ los elementos de
$\mathcal{F}(U)$, y se escribe entonces $\Gamma(U,\mathcal{F})$ en lugar de
$\mathcal{F}(U)$; para $s\in\Gamma(U,\mathcal{F})$ y $V$ abierto contenido en
$U$, se escribe $s|V$ en lugar de $\rho_V^U(s)$; para todo $x\in U$, la imagen
canónica de $s$ en $\mathcal{F}_x$ es el \emph{germen} de $s$ en el punto $x$,
denotado por $s_x$ (\emph{nunca emplearemos la notación $s(x)$ en este sentido},
pues esta notación se reserva para otra noción relativa a los haces
particulares que se considerarán en este Tratado (5.5.1)).

Si entonces $u:\mathcal{F}\to\mathcal{G}$ es un morfismo de haces con valores
en $\boldsymbol{K}$, se escribirá $u(s)$ en lugar de $u_V(s)$ para todo
$s\in\Gamma(U,\mathcal{F})$.

Si $\mathcal{F}$ es un haz de grupos conmutativos, de anillos o de
módulos, se dice que el conjunto de los $x\in X$ tales que
$\mathcal{F}_x\neq\{0\}$ es el \emph{soporte} de $\mathcal{F}$, denotado por
$\operatorname{Supp}(\mathcal{F})$; este conjunto no es necesariamente
cerrado en $X$.

Cuando $\boldsymbol{K}$ está definida por una especie de estructura con morfismos,
\emph{nos abstendremos sistemáticamente de recurrir al punto de
vista de los «espacios étalé»} en lo relativo a los haces con valores en
$\boldsymbol{K}$; es decir, nunca consideraremos un haz como
un espacio topológico (ni siquiera como el conjunto unión de sus tallos), y
tampoco consideraremos un morfismo
$u:\mathcal{F}\to\mathcal{G}$ de tales haces sobre $X$ como una aplicación
continua de espacios topológicos.
\end{env}
